% GENERATED FILE. Edit canonical lessons, then run npm run generate:books.
% canonical-source-hash: fnv1a64:d8c4a2361f3045bd
% canonical-lessons: RU-W06-u, RU-W06-z, RU-W06-k, RU-W06-l, RU-W06-m, RU-W06-y, RU-W06-zh, RU-W06-ch, RU-W06-sh, RU-W06-soft-sign, RU-W06-ya

\chapter{Eleven Letters, Sorted Into Three Kinds}
\label{ch:ru-script-eleven}
\hlchaptermodality{14}

\begin{chapteropening}
\textbf{By the end of this chapter:} \emph{I can write eleven more Cyrillic letters from memory, sort any Cyrillic letter into true friend, false friend or new shape, and explain what the soft sign does even though it makes no sound of its own.}

A letter that is a whole syllable and looks like a backwards R that it has nothing to do with — closing a chapter whose organising idea is that every Cyrillic letter sorts into one of three kinds.
\end{chapteropening}

\section[u]{\ru{у} — the letter that looks like Latin y and says 'oo'}

\label{lesson:RU-W06-u}



\begin{quote}
You have met the false friends and the first new shapes. This chapter takes the eleven letters that block the most words you cannot yet read.
\end{quote}

\subsection*{Script}
\begin{quote}
\ru{у}
\end{quote}

It looks like a Latin \textbf{y}. It is not one. It says \textbf{oo}, the vowel of English \emph{boot}.

\textbf{This is the single most expensive letter to get wrong}, because it is common: it blocks more words in this book than any other letter you have left. Every time an English eye reads it as \emph{y}, a word comes apart.

The shape is close to Latin \emph{y} — a small bowl on top, a tail below the line — but the tail is usually straighter and the bowl smaller.

Say it round and short, with the lips pushed forward. Not the \emph{y} of \emph{yes}. Not the \emph{u} of \emph{cup}. The \emph{oo} of \emph{boot}.

\subsection*{Your turn}
\begin{itemize}
\raggedright
  \item \emph{Write it:} \ru{у}
  \item \emph{Sort:} true friend, false friend, or new shape?
\end{itemize}

\begin{tcolorbox}[breakable,colback=teal!4,colframe=teal!35!black,title={Before you move on}]
What does \textbf{\ru{у}} sound like? (\textbf{oo}, as in \emph{boot}.) What does an English eye want it to be? (A \textbf{y} — and that is the trap.)

Source: \href{https://www.unicode.org/charts/PDF/U0400.pdf}{Unicode Cyrillic chart}.
\end{tcolorbox}
\section[z]{\ru{з} — a letter shaped like the digit 3}

\label{lesson:RU-W06-z}



\begin{quote}
Write the letter from the previous lesson, and say which of the three kinds it is.
\end{quote}

\subsection*{Script}
\begin{quote}
\ru{з}
\end{quote}

\textbf{z}, the \emph{z} of English \emph{zoo}. And the shape is exactly what it looks like: \textbf{the digit 3}.

That is not a coincidence to memorise around — it is simply a convenient hook, and a rare piece of luck in this alphabet. Nothing in Latin looks like this, so nothing can mislead you.

Write it as you would write a 3, and read it as \emph{z}.

\subsection*{Your turn}
\begin{itemize}
\raggedright
  \item \emph{Write it:} \ru{з}
  \item \emph{Sort:} true friend, false friend, or new shape?
\end{itemize}

\begin{tcolorbox}[breakable,colback=teal!4,colframe=teal!35!black,title={Before you move on}]
What does \textbf{\ru{з}} sound like? (\textbf{z}.) And what does it look like? (The \textbf{digit 3}.)

Source: \href{https://www.unicode.org/charts/PDF/U0400.pdf}{Unicode Cyrillic chart}.
\end{tcolorbox}
\section[k]{\ru{к} — a true friend}

\label{lesson:RU-W06-k}



\begin{quote}
Write the letter from the previous lesson, and say which of the three kinds it is.
\end{quote}

\subsection*{Script}
\begin{quote}
\ru{к}
\end{quote}

\textbf{k}, the \emph{k} of English \emph{skip}. It looks like a Latin \textbf{k} and it \textbf{is} one.

Cyrillic letters come in three kinds for an English reader, and it is worth having the three named:

\begin{itemize}
\raggedright
  \item \textbf{true friends} — look Latin, sound Latin. This one, and the letters romanised \emph{a}, \emph{o}, \emph{m}, \emph{t}.
  \item \textbf{false friends} — look Latin, sound different. \textbf{\ru{у}} you just met, and the ones you met earlier as \emph{v}, \emph{r}, \emph{s} and \emph{n}.
  \item \textbf{new shapes} — look like nothing Latin at all. \textbf{\ru{з}} from a moment ago, and the two you learned as \emph{b} and \emph{d}.
\end{itemize}

There is no fourth kind. Sorting every new letter into one of those three is most of what learning this alphabet is.

\subsection*{Your turn}
\begin{itemize}
\raggedright
  \item \emph{Write it:} \ru{к}
  \item \emph{Sort:} true friend, false friend, or new shape?
\end{itemize}

\begin{tcolorbox}[breakable,colback=teal!4,colframe=teal!35!black,title={Before you move on}]
What are the three kinds of Cyrillic letter for an English reader? (\textbf{True friends}, \textbf{false friends}, \textbf{new shapes}.) Which is \textbf{\ru{к}}? (A \textbf{true friend}.)

Source: \href{https://www.unicode.org/charts/PDF/U0400.pdf}{Unicode Cyrillic chart}.
\end{tcolorbox}
\section[l]{\ru{л} — a new shape}

\label{lesson:RU-W06-l}



\begin{quote}
Write the letter from the previous lesson, and say which of the three kinds it is.
\end{quote}

\subsection*{Script}
\begin{quote}
\ru{л}
\end{quote}

\textbf{l}, the \emph{l} of English \emph{love}. A \textbf{new shape}: nothing in Latin looks like it.

Two strokes meeting at a peak, with the left leg kicking out. Written quickly it leans, and in some fonts it is nearly a triangle with an open base.

Careful with one thing. In handwriting \textbf{\ru{л}} and \textbf{\ru{м}} are close relatives, and you are about to meet \textbf{\ru{м}}. The difference is the top: \textbf{\ru{л}} comes to a \textbf{point}, \textbf{\ru{м}} has a \textbf{dip} in the middle.

\subsection*{Your turn}
\begin{itemize}
\raggedright
  \item \emph{Write it:} \ru{л}
  \item \emph{Sort:} true friend, false friend, or new shape?
\end{itemize}

\begin{tcolorbox}[breakable,colback=teal!4,colframe=teal!35!black,title={Before you move on}]
Is \textbf{\ru{л}} a true friend, a false friend, or a new shape? (A \textbf{new shape}.) How do you tell it from the letter coming next? (\textbf{\ru{л}} peaks; the other \textbf{dips}.)

Source: \href{https://www.unicode.org/charts/PDF/U0400.pdf}{Unicode Cyrillic chart}.
\end{tcolorbox}
\section[m]{\ru{м} — a true friend, and its capital}

\label{lesson:RU-W06-m}



\begin{quote}
Write the letter from the previous lesson, and say which of the three kinds it is.
\end{quote}

\subsection*{Script}
\begin{quote}
\ru{м}   \ru{М}
\end{quote}

\textbf{m}, the \emph{m} of English \emph{met} — a \textbf{true friend} in both cases.

And here is a fact about this alphabet worth having early: \textbf{most Cyrillic capitals are simply the lowercase letter drawn large.} Latin makes you learn two shapes for \emph{a}, \emph{b}, \emph{d}, \emph{e}, \emph{g}; Cyrillic mostly does not.

This letter is one of them, and so are the ones romanised \emph{k}, \emph{o}, \emph{s} and \emph{t}: same shape, two sizes.

\textbf{Not all of them} — the letters for \emph{b}, \emph{d} and \emph{a} do have distinct capitals — but the default is that a capital costs you nothing new, which is a genuine saving over the Latin alphabet and the opposite of what a beginner expects.

\subsection*{Your turn}
\begin{itemize}
\raggedright
  \item \emph{Write it:} \ru{м}
  \item \emph{Sort:} true friend, false friend, or new shape?
\end{itemize}

\begin{tcolorbox}[breakable,colback=teal!4,colframe=teal!35!black,title={Before you move on}]
What sound is \textbf{\ru{м}}? (\textbf{m}.) And what is true of most Cyrillic capitals? (They are the \textbf{lowercase drawn large} — no new shape to learn.)

Source: \href{https://www.unicode.org/charts/PDF/U0400.pdf}{Unicode Cyrillic chart}.
\end{tcolorbox}
\section[y]{\ru{ы} — the vowel English has no name for}

\label{lesson:RU-W06-y}



\begin{quote}
Write the letter from the previous lesson, and say which of the three kinds it is.
\end{quote}

\subsection*{Script}
\begin{quote}
\ru{ы}
\end{quote}

Now a letter with no English equivalent at all — not a shape problem, a \textbf{sound} problem.

\textbf{\ru{ы}} is a vowel made with the tongue pulled \textbf{back and high}, roughly where you hold it for \emph{oo}, but with the lips \textbf{unrounded} and spread as for \emph{ee}. English has no such vowel and no letter for it. Books usually transliterate it \emph{y}, which helps nobody.

The reliable way in: say \emph{ee}. Hold it. Now, without moving your lips, pull your tongue back towards your throat. What comes out is close.

The shape is two pieces — a hard sign body and a vertical — and it is \textbf{never capitalised at the start of a word}, because no Russian word begins with it.

\subsection*{Your turn}
\begin{itemize}
\raggedright
  \item \emph{Write it:} \ru{ы}
  \item \emph{Sort:} true friend, false friend, or new shape?
\end{itemize}

\begin{tcolorbox}[breakable,colback=teal!4,colframe=teal!35!black,title={Before you move on}]
Is \textbf{\ru{ы}} a shape problem or a sound problem? (A \textbf{sound} problem.) How do you find it? (Say \emph{ee}, then \textbf{pull the tongue back} without moving the lips.) And can a word begin with it? (\textbf{No}.)

Source: \href{https://www.unicode.org/charts/PDF/U0400.pdf}{Unicode Cyrillic chart}.
\end{tcolorbox}
\section[zh]{\ru{ж} — the beetle}

\label{lesson:RU-W06-zh}



\begin{quote}
Write the letter from the previous lesson, and say which of the three kinds it is.
\end{quote}

\subsection*{Script}
\begin{quote}
\ru{ж}
\end{quote}

\textbf{zh}, the sound in the middle of English \emph{measure} or \emph{vision}. A \textbf{new shape}, and the most recognisable one in Cyrillic: three strokes radiating from a central spine, symmetrical, often called \textbf{the beetle}.

The sound is worth isolating. English has it — \emph{measure}, \emph{treasure}, \emph{azure}, \emph{genre} — but spells it with no letter of its own, and no ordinary English word opens with it. Russian gives it a letter and opens words with it freely.

Symmetrical, so there is no wrong way round. Draw the spine, then the four legs.

\subsection*{Your turn}
\begin{itemize}
\raggedright
  \item \emph{Write it:} \ru{ж}
  \item \emph{Sort:} true friend, false friend, or new shape?
\end{itemize}

\begin{tcolorbox}[breakable,colback=teal!4,colframe=teal!35!black,title={Before you move on}]
What sound is \textbf{\ru{ж}}? (The middle of \emph{measure}.) Does English have that sound? (\textbf{Yes} — but no letter for it.) And what is the shape called? (\textbf{The beetle}.)

Source: \href{https://www.unicode.org/charts/PDF/U0400.pdf}{Unicode Cyrillic chart}.
\end{tcolorbox}
\section[ch]{\ru{ч} — one letter for a sound English spells with two}

\label{lesson:RU-W06-ch}



\begin{quote}
Write the letter from the previous lesson, and say which of the three kinds it is.
\end{quote}

\subsection*{Script}
\begin{quote}
\ru{ч}
\end{quote}

\textbf{ch}, the \emph{ch} of English \emph{cheese}. A \textbf{new shape} — roughly a Latin \emph{h} rotated, or a wine glass.

Notice what the alphabet is doing. English needs \textbf{two} letters, \emph{c} and \emph{h}, to write this sound, and that pair is ambiguous — compare \emph{church}, \emph{chorus}, \emph{machine}. Cyrillic spends \textbf{one letter} on it and the letter means one thing.

This is a genuine advantage of the script and it is worth naming: Russian spelling is far more regular than English spelling. A letter usually says one thing.

\subsection*{Your turn}
\begin{itemize}
\raggedright
  \item \emph{Write it:} \ru{ч}
  \item \emph{Sort:} true friend, false friend, or new shape?
\end{itemize}

\begin{tcolorbox}[breakable,colback=teal!4,colframe=teal!35!black,title={Before you move on}]
How many letters does English use for this sound? (\textbf{Two} — and ambiguously.) How many does Russian? (\textbf{One}.)

Source: \href{https://www.unicode.org/charts/PDF/U0400.pdf}{Unicode Cyrillic chart}.
\end{tcolorbox}
\section[sh]{\ru{ш} — the comb}

\label{lesson:RU-W06-sh}



\begin{quote}
Write the letter from the previous lesson, and say which of the three kinds it is.
\end{quote}

\subsection*{Script}
\begin{quote}
\ru{ш}
\end{quote}

\textbf{sh}, the \emph{sh} of English \emph{shoe}. A \textbf{new shape}: three vertical teeth on a base, like a \textbf{comb}.

Same point as the last letter, and it is a pattern rather than a coincidence: English writes this with two letters, Russian with one.

Keep it apart from the beetle by counting: \textbf{\ru{ж}} radiates from a centre, \textbf{\ru{ш}} stands on a flat base with three uprights.

\subsection*{Your turn}
\begin{itemize}
\raggedright
  \item \emph{Write it:} \ru{ш}
  \item \emph{Sort:} true friend, false friend, or new shape?
\end{itemize}

\begin{tcolorbox}[breakable,colback=teal!4,colframe=teal!35!black,title={Before you move on}]
What sound is \textbf{\ru{ш}}? (\textbf{sh}.) How do you tell it from the beetle? (\textbf{\ru{ш}} has three uprights on a \textbf{flat base}; the beetle \textbf{radiates} from a centre.)

Source: \href{https://www.unicode.org/charts/PDF/U0400.pdf}{Unicode Cyrillic chart}.
\end{tcolorbox}
\section[(soft sign)]{\ru{ь} — a letter that makes no sound of its own}

\label{lesson:RU-W06-soft-sign}



\begin{quote}
Write the letter from the previous lesson, and say which of the three kinds it is.
\end{quote}

\subsection*{Script}
Now the idea this alphabet has and Latin does not.

\begin{quote}
\ru{ь}
\end{quote}

This is the \textbf{soft sign}, and \textbf{it is not a sound}. You cannot pronounce it. It has no name you say aloud in a word. What it does is change the letter \textbf{before} it: it \textbf{softens} that consonant, which means the tongue moves towards the roof of the mouth, as if a faint \emph{y} were sliding in after it.

English does this by accident. Compare the \emph{n} in \emph{canyon} with the \emph{n} in \emph{can} — the first is softer, because a \emph{y} follows. Russian treats that difference as \textbf{meaningful} and gives it a letter.

\textbf{A letter that is a modifier rather than a sound} is the one genuinely alien thing in this script for an English reader. Everything else is a shape to learn or a sound to hear; this is a different kind of object entirely.

It is also common — it blocks a great many words in this book — and it never begins a word, so you will only ever meet it after a consonant.

\subsection*{Your turn}
\begin{itemize}
\raggedright
  \item \emph{Write it:} \ru{ь}
  \item \emph{Sort:} true friend, false friend, or new shape?
\end{itemize}

\begin{tcolorbox}[breakable,colback=teal!4,colframe=teal!35!black,title={Before you move on}]
Can you pronounce \textbf{\ru{ь}}? (\textbf{No} — it makes no sound of its own.) What does it do? (\textbf{Softens} the consonant before it.) And where in a word can it never appear? (At the \textbf{start}.)

Source: \href{https://www.unicode.org/charts/PDF/U0400.pdf}{Unicode Cyrillic chart}.
\end{tcolorbox}
\section[ya]{\ru{Я} — one letter that is a whole syllable}

\label{lesson:RU-W06-ya}



\begin{quote}
Write the letter from the previous lesson, and say which of the three kinds it is.
\end{quote}

\subsection*{Script}
\begin{quote}
\ru{Я}   \ru{я}
\end{quote}

Last, and it carries two lessons at once.

\textbf{It looks like a backwards Latin R.} It is not an R and has nothing to do with one. That resemblance is the reason this letter appears on every parody of Russian ever printed in English, and it means nothing.

\textbf{And it is not one sound but two.} \textbf{\ru{я}} says \textbf{ya} — a whole syllable, a \emph{y} glide plus an \emph{a}, in a single letter. Cyrillic has several of these, and they are why a Russian word can look shorter than its pronunciation.

On its own, \textbf{\ru{я}} is also the Russian word for \textbf{I}. A single letter, a whole word — like English \emph{a} and \emph{I}, and for the same reason: the commonest words wear down to the shortest forms.

That is your eleventh piece. Sort each one into the three kinds as you go, and the rest of the alphabet is more of the same.

\subsection*{Your turn}
\begin{itemize}
\raggedright
  \item \emph{Write it:} \ru{Я}
  \item \emph{Sort:} true friend, false friend, or new shape?
\end{itemize}

\begin{tcolorbox}[breakable,colback=teal!4,colframe=teal!35!black,title={Before you move on}]
What does \textbf{\ru{я}} sound like? (\textbf{ya} — two sounds, one letter.) Is it related to Latin R? (\textbf{No} — the resemblance is meaningless.) And what word is it on its own? (\textbf{I}.)

Source: \href{https://www.unicode.org/charts/PDF/U0400.pdf}{Unicode Cyrillic chart}.
\end{tcolorbox}
